\chapter{Контракт Stage 1: минимальный срез, который строим первым}
\label{ch:stage1}

Монография зрелая, и улучшать её можно бесконечно. Но полезнее не ещё одна красивая
глава, а \emph{зафиксированный маленький вертикальный срез}, который можно собрать за
неделю-две и начать \emph{ломать} на стенде. Это приложение — контракт такой сборки:
не весь организм памяти, а сквозная линия от захвата до речи.

\section{Сквозной срез (pipeline v0)}

\begin{lstlisting}[style=pseudo]
event_log            # неизменяемый лог (гл. 17): append observed/committed/deleted
  -> typed nodes     # схема узла v0 (ниже), node_kind обязателен
  -> write/skip gate # epistemic + salience + anti-fusion; skip логируется (гл. 17)
  -> forensic index  # точный поиск по метаданным/провенансу/времени (гл. 17)
  -> anchor retrieval# вектор-якоря + ограниченный обход (гл. 6)
  -> context-packer v0 # с метками observed/hypothesis/low_conf/... (гл. 17)
  -> inspector CLI   # показать узлы/рёбра/поля/лог, вручную инспектировать
  -> synthetic tests # маленький граф с ground truth (гл. 18)
\end{lstlisting}

\section{Схема узла v0 (минимум полей)}

\texttt{id}, \texttt{content}, \texttt{embedding}, \texttt{node\_kind},
\texttt{confidence}$\in[-1,1]$, \texttt{provenance} (или support-skeleton),
\texttt{time}, \texttt{layer}. Остальные поля (\texttt{valence},
\texttt{importance}, \texttt{domains}, \texttt{need\_to\_*}) добавляются по мере
надобности, но \texttt{node\_kind} обязателен с самого начала — он задаёт смысл
\texttt{confidence} (гл.~\ref{ch:substrate}).

\section{Первые 6--8 сущностей (что реально кодим)}

\begin{enumerate}\itemsep1pt
\item \textbf{Event log} (append-only) $+$ пересборка графа-проекции из лога.
\item \textbf{Typed node store} (KV) $+$ edge store (adjacency).
\item \textbf{Write/skip gate} с логированием причины skip.
\item \textbf{Forensic index} (точный поиск/удаление по метаданным) — включая
работающий delete-cascade (гл.~\ref{ch:runtime}, раздел~\ref{sec:forget}).
\item \textbf{Anchor retrieval} (вектор-индекс $+$ ограниченный обход).
\item \textbf{Context-packer v0} (с эпистемическими метками).
\item \textbf{Inspector CLI} — «memory inspector»: смотреть узлы, рёбра, поля, лог.
\item \textbf{Synthetic graph tests} (SBM/BA $+$ подсаженные сценарии).
\end{enumerate}

\section{Что сознательно игнорируем на Stage 1}

Чтобы срез остался тонким: mood-динамика и её слои, личность/Rinness и drive-вектор,
сон/скраб и разрешение противоречий, тяжёлый fusion (только anti-fusion-скелет и
провизорные кандидаты), гиперрёбра и клеточные комплексы, governance-слой,
мета-обучение операторов. Это не отменяется — это \emph{откладывается}; красивые
главы ждут, пока нижний срез не станет доверенным.

Три конкретных инженерных упрощения на первую итерацию (полные механизмы описаны, но
дороги в реализации):
\begin{itemize}\itemsep2pt
\item \textbf{Удаление — только \texttt{tombstone}.} Crypto-shredding и
taint-propagation (гл.~\ref{ch:runtime}) безупречны, но каскад «яда» по кэшам, индексам
и бэкапам съест месяцы. В прототипе — tombstone $+$ пересборка проекции; криптостирание
позже.
\item \textbf{Коллизии консолидации — компоненты связности, не модулярность.} Полный
со-консолидационный граф $+$ community detection тяжёл. На старте: простые
\emph{connected components} по графу предложений с отсечением по жёсткому порогу
пересечения (LSH\,/\,MinHash). Модулярность — когда упрёмся.
\item \textbf{Семантический вентиль — только косинус кэш-эмбеддингов.}
\texttt{sense/type}-гейт против омонимии (гл.~\ref{ch:walker}) требует sense-слоя,
которого ещё нет. \emph{Implementation risk}: Stage 1 \emph{сознательно не обещает}
сильную защиту от омонимии — либо терпим протечки, либо руками делаем крошечную
sense-схему на 20--30 демо-кейсов. Лечим, когда эмпирически упрёмся.
\end{itemize}

\section{Критерии готовности v0 (как поймём, что работает)}

\begin{itemize}\itemsep1pt
\item запись и \emph{skip} работают, причина skip видна в логе;
\item \texttt{tombstone}-удаление сносит узел из проекции, это проверяемо inspector-ом
(taint-каскад — отдельная поздняя цель);
\item anchor-retrieval на синтетическом графе возвращает осмысленную окрестность;
\item реплей event-log \emph{перестраивает} тот же граф-проекцию (event sourcing);
\item хотя бы один тест из батареи отказов (гл.~\ref{ch:experiment}) — например,
$\pi$-концентрация на хабах или точное удаление — зелёный.
\end{itemize}

\begin{principle}[Первый успех — не «Рин помнит», а «физика графа наблюдаема»]
Формулировка критерия важнее списка. Если первым мерилом сделать «Рин умно пользуется
памятью в диалоге», сразу утонешь в LLM-шуме и не поймёшь, работает граф или врёт
модель. Поэтому Stage 1 — это \emph{memory inspector v0}, а не «память v0»: успех —
когда мы можем руками создать event, увидеть узел, увидеть retrieval, увидеть proposal,
увидеть skip\,/\,commit, увидеть support-skeleton. Наблюдаемость внутренней механики
раньше, чем «умное поведение».
\end{principle}

\begin{principle}[Дальше — строить и ломать, а не расширять]
На этом теоретический слой монографии \emph{заморожен}. Следующий лучший ход —
собрать этот срез, прогнать synthetic tests, и позволить реальности править
гипотезы. Каждая «красивая глава» включается только тогда, когда нижний срез её
выдержит.
\end{principle}
